\documentclass{article}

\usepackage{amssymb,amsmath,mleftright}

\begin{document}

General model assumptions:
\begin{itemize}
\item the fluid density and viscosity are constant in time and space (implies incompressibility);
\item the fluid flow is laminar;
\item the column is operated under constant conditions (e.g., temperature, flow rate);
\item the process is isothermal (i.e., there are no thermal effects);
\item diffusion does not depend on the concentration of the components, viscosity of the fluid, or pressure;
\item the fluid in the particles is stagnant (i.e., there is no convective flow);
\item the porous particles are spherical, rigid, and do not move;
\item the particles are homogeneous and of uniform porosity;
\item all components access the same particle volume;
\item the partial molar volumes are the same in mobile and solid phase;
\item the binding model parameters do not depend on pressure and are constant along the column;
\item the solvent is not adsorbed;
\end{itemize}


Specific model assumptions:
\begin{itemize}
\item the fluid only flows in the axial direction of the column (i.e., there is no flow in the radial and angular direction);
\item the column has a conical frustum shape with linearly varying radius;
\item the particles form a continuum inside the column (i.e., there is interstitial and particle volume at every point in the column);
\item the column is radially symmetric and homogeneous (i.e., concentration profiles and parameters only depend on the axial position);
\item the interstitial liquid phase concentration is spatially constant on the outer particle surface;
\item the pore and surface diffusion are infinitely fast. That is, we assume $D_{i}^{\mathrm{p}} = D_{i}^{\mathrm{s}} = \infty$;
\item the film around the particles does not limit mass transfer. That is, we assume $k^{\mathrm{f}}_i = \infty$)
\end{itemize}


\section*{Frustum Lumped Rate Model without Pores}
Consider a conical frustum column of length $L > 0$ with inlet radius $R^0 > 0$ and outlet radius $R^L > 0$ packed with spherical particles, and observed over a time interval $(0, T^{\mathrm{end}})$.
In the interstitial volume, mass transfer is governed by the following convection-diffusion equations in $(0, T^\mathrm{end})\times (0, L)$ and for all components $i\in\{1, \dots, N^{\mathrm{c}} \}$
\begin{align}
\frac{\partial c^{\mathrm{b}}_i}{\partial t} + \frac{1 - \varepsilon^{\mathrm{t}}}{\varepsilon^{\mathrm{t}}} \frac{\partial c^{\mathrm{s}}_i}{\partial t} = - \frac{v}{r(x)^2} \frac{\partial c^{\mathrm{b}}_i}{\partial x} + \frac{1}{r(x)^2} \frac{\partial}{\partial x} \left( r(x)^2 \varepsilon^{\mathrm{c}} D^{\mathrm{ax}}_{i} \frac{\partial c^{\mathrm{b}}_i}{\partial x} \right),
\end{align}
with boundary conditions

\begin{alignat}{2}
\frac{v}{r(x)^2} c_{\mathrm{in},i} &= \left.\left( \frac{v}{r(x)^2} c^{\mathrm{b}}_i - D^{\mathrm{ax}}_{i} \frac{\partial c^{\mathrm{b}}_i}{\partial x}\right)\right|_{x=0} & &\qquad\text{on }(0, T^{\mathrm{end}}),\\
               0 &= - D^{\mathrm{ax}}_{i} \left. \frac{\partial c^{\mathrm{b}}_i}{\partial x} \right|_{x=L} & &\qquad\text{on }(0, T^{\mathrm{end}}).
\end{alignat}
Here, $t\in (0, T^{\mathrm{end}})$ is the time coordinate, $x\in (0, L)$ is the axial coordinate, $c^{\mathrm{b}}_i\colon (0, T^\mathrm{end})\times (0, L) \mapsto \mathbb{R}$ is the bulk liquid concentration, $Q> 0$ is the volumetric flow rate, $R(x) = R^0 + \frac{R^L - R^0}{L} x$ is the column radius function, $\varepsilon^{\mathrm{t}}\in (0, 1)$ is the total porosity, $v:= \frac{Q}{\pi}$ is the velocity coefficient, and $\tilde{D}^\mathrm{ax}_i\geq 0$ is the apparent axial dispersion coefficient.
In the particles, mass transfer is governed by reaction equations in $(0, T^\mathrm{end}) \times (0, L)$ and for all components
\begin{align}
          \frac{\partial c^{\mathrm{s}}_{i}}{\partial t}
          &= f^{\mathrm{bind}}_{i}\left( \vec{c}^{\mathrm{\ell}}, \vec{c}^{\mathrm{s}} \right), \end{align}
Here, $c^{\mathrm{s}}_{i}\colon (0, T^\mathrm{end}) \times (0, L) \mapsto \mathbb{R}$ is the particle solid concentration, $f^\mathrm{bind}_{i}$ is the adsorption isotherm function, $\vec{c}^\mathrm{p}$ is the particle liquid components vector, and $\vec{c}^\mathrm{s}$ is the particle solid components vector.
Consistent initial values for all solution variables (concentrations) are defined at $t = 0$.
\end{document}